% Part: second-order-logic
% Chapter: syntax-and-semantics
% Section: semantic-notions

\documentclass[../../../include/open-logic-section]{subfiles}

\begin{document}

\olfileid{sol}{syn}{sem}
\olsection{معنايي مفهومونه}

\begin{explain}
د \emph{اعتبار}، \emph{استلزام} او \emph{د صدق وړتيا} مرکزي منطقي
مفهومونه د دويمې درجې منطق دپاره د لومړۍ درجې منطق په شان تعريفېږي؛
يوازې بنسټيزه د صدق اړيکه اوس د دويمې درجې د !!{formula}s اړيکه ده۔
د دويمې درجې !!{sentence} هغه !!a{formula} ده چې ټول متغيرونه
يې، د محمول او تابع متغيرونو په ګډون، تړلي وي۔
\end{explain}

\begin{defn}[اعتبار]
جمله~$!A$ هله \emph{معتبره} ده، $\Entails !A$، چې د هر
!!{structure}~$\Struct M$ دپاره $\Sat{M}{!A}$۔
\end{defn}

\begin{defn}[استلزام]
د جملو سټ~$\Gamma$ هله جمله~$!A$ \emph{لازموي}، $\Gamma
\Entails !A$، چې د هر هغه !!{structure}~$\Struct M$ دپاره چې
$\Sat{M}{\Gamma}$ وي، $\Sat{M}{!A}$ هم وي۔
\end{defn}

\begin{defn}[د صدق وړتيا]
د جملو سټ~$\Gamma$ هله \emph{د صدق وړ} دے چې د کوم
!!{structure}~$\Struct M$ دپاره $\Sat{M}{\Gamma}$ وي۔ که~$\Gamma$
د صدق وړ نۀ وي، \emph{د صدق نۀ وړ} بلل کېږي۔
\end{defn}

\end{document}
